\chapter{Rules}\label{ch:rules}
\chapterintro{Thirteen deterministic rules check that a specification states an admissible problem before anything is
solved. This chapter explains when the rules run, how to read a finding, and what each code means and how to fix it.
The rule list in this chapter is generated from \code{src/agentto/rules.py}.}

\section{When the rules run}
\index{rules!when they run}
\begin{itemize}
  \item In the workspace, \CheckDelay~ms after the last edit. In the paper's latency study the 95th-percentile
  validation time per task was 1.2--10.5~ms \claim{C6.4}.
  \item On every \tool{validate_spec} call, at the specification's own mesh or at a mesh size \code{h} you pass.
  \item On every \tool{solve} call, again at the mesh of the requested rung. A solve whose rung fails validation is
  refused before it reaches the GPU, and the GPU worker checks once more before it starts.
\end{itemize}

Any \sev{error} blocks solving and release. A \sev{warning} or \sev{info} finding does not block; it stays visible in
the \gui{Checks} panel, often with a suggested fix. Several rules depend on the grid (for example R002, R010, R011 and
R013), so a specification can pass at one mesh and fail at another.

\section{Reading a finding}
Each finding has a stable code, a severity, a message, the ids of the items it concerns (the tree marks those items),
and, where one exists, a fix. A fix either sets a field (the workspace can offer it as one click, and an agent can
propose it as a card) or suggests an action in words. \tool{validate_spec} returns the findings as JSON:

\begin{jsoncode}
{"ok": false,
 "resolution": {"h": 1.0, "grid": [100, 50, 50], "n_elem": 250000, ...},
 "findings": [{"code": "R003", "severity": "error",
   "message": "supports leave rigid-body motion free: translation y, translation z, rotation about x",
   "ids": ["support_1"],
   "fix": {"suggest": "add a support that restrains translation y, translation z, rotation about x"}}]}
\end{jsoncode}

The codes are stable: experiments, logs and this book refer to them, and a code is never reused for a different
check.

\section{Summary}
% Generated by docs/guidebook/gen_tables.py from the code. Do not edit by hand.
\begin{reftable}{@{}l>{\raggedright\arraybackslash}p{0.2\linewidth}X@{}}
Code & Severity & Fires when \\ \midrule
\hyperref[rule:R001]{\textbf{R001}} & \sev{error} & no grid fits the domain at this element size \\
\hyperref[rule:R002]{\textbf{R002}} & \sev{error} & a support selects no nodes \\
\hyperref[rule:R003]{\textbf{R003}} & \sev{error} & the supports leave a rigid-body mode free \\
\hyperref[rule:R004]{\textbf{R004}} & \sev{error} & a load selects no nodes, or every load is zero \\
\hyperref[rule:R005]{\textbf{R005}} & \sev{error} / \sev{warning} & part of the load acts on fixed degrees of freedom \\
\hyperref[rule:R006]{\textbf{R006}} & \sev{error} & no design elements remain \\
\hyperref[rule:R007]{\textbf{R007}} & \sev{error} & elements are both keep-in and keep-out \\
\hyperref[rule:R008]{\textbf{R008}} & \sev{error} & a support or load lies entirely inside a keep-out region \\
\hyperref[rule:R009]{\textbf{R009}} & \sev{warning} & the minimum feature is not resolvable at this mesh \\
\hyperref[rule:R010]{\textbf{R010}} & \sev{warning} & a region was snapped to the grid \\
\hyperref[rule:R011]{\textbf{R011}} & \sev{warning} & a passive region contains no element centre \\
\hyperref[rule:R012]{\textbf{R012}} & \sev{warning} / \sev{info} & supports change on the coarsest multigrid level \\
\hyperref[rule:R013]{\textbf{R013}} & \sev{error} & the grid needs more GPU memory than is available \\
\end{reftable}


\section{Rule by rule}
In the messages below, \meta{id} and similar placeholders stand for values filled in at run time.

% Generated by docs/guidebook/gen_tables.py from the code. Do not edit by hand.
\begin{rulecard}{R001}{\sev{error}}{no grid fits the domain at this element size}\label{rule:R001}\index{R001@\texttt{R001} (rule)}
\rulefield{Message}\msgline{domain \meta{axis} size \meta{length} is not a whole number of elements of edge \meta{h}}
\rulefield{Meaning}A side of the domain is not a whole number of elements of edge $h$ (or the number of cells along the short edge gives a non-integer count on another side). R001 stops the check: no other rule runs until a grid exists.
\rulefield{Fix}Choose an $h$ that divides all three sides, or accept the one-click fix, which sets \code{cells\_short\_edge} to a value that fits. \emph{(The finding carries a fix: one click: set \code{resolution.cells\_\allowbreak{}short\_\allowbreak{}edge}.)}
\end{rulecard}

\begin{rulecard}{R002}{\sev{error}}{a support selects no nodes}\label{rule:R002}\index{R002@\texttt{R002} (rule)}
\rulefield{Message}\msgline{support '\meta{id}' selects no nodes}
\rulefield{Meaning}The support's region misses every grid node: a box between grid planes, a cylinder thinner than the node spacing, an edge span outside the domain. The support does nothing.
\rulefield{Fix}Run \code{list\_regions} to see the node count of each item, then move or enlarge the region, or refine the mesh.
\end{rulecard}

\begin{rulecard}{R003}{\sev{error}}{the supports leave a rigid-body mode free}\label{rule:R003}\index{R003@\texttt{R003} (rule)}
\rulefield{Message}\msgline{supports leave rigid-body motion free: \meta{modes}}
\rulefield{Meaning}The fixed degrees of freedom do not remove all six rigid-body modes, so the stiffness matrix is singular. The message names the free translations and rotations. A common cause is a single edge support (free rotation about the edge) or rollers that fix only one direction.
\rulefield{Fix}Add a support, or fix more directions, so that the listed modes are restrained. Fix in the spec; do not rely on the solver to fail. \emph{(The finding carries a fix: suggests: add a support that restrains \meta{modes}.)}
\end{rulecard}

\begin{rulecard}{R004}{\sev{error}}{a load selects no nodes, or every load is zero}\label{rule:R004}\index{R004@\texttt{R004} (rule)}
\rulefield{Message}\msgline{every load has zero force} \newline \msgline{load '\meta{id}' selects no nodes}
\rulefield{Meaning}Either a load region misses every node, or the force vectors of all loads are zero, so there is nothing to optimize against.
\rulefield{Fix}Check the load's region with \code{list\_regions} and its force components with \code{explain\_spec}.
\end{rulecard}

\begin{rulecard}{R005}{\sev{error} / \sev{warning}}{part of the load acts on fixed degrees of freedom}\label{rule:R005}\index{R005@\texttt{R005} (rule)}
\rulefield{Message}\msgline{\meta{share}\% of the applied load acts on fixed DOFs and is lost}
\rulefield{Meaning}The problem builder drops every load component that falls on a fixed degree of freedom (it would go straight into the support) and records the share. Above 50\% of the load the finding is an error, otherwise a warning.
\rulefield{Fix}Move the load off the support, or leave that direction free at the shared nodes.
\end{rulecard}

\begin{rulecard}{R006}{\sev{error}}{no design elements remain}\label{rule:R006}\index{R006@\texttt{R006} (rule)}
\rulefield{Message}\msgline{no design elements remain after passive regions}
\rulefield{Meaning}The keep-in and keep-out regions together cover every element, so there is nothing to design.
\rulefield{Fix}Shrink or remove passive regions.
\end{rulecard}

\begin{rulecard}{R007}{\sev{error}}{elements are both keep-in and keep-out}\label{rule:R007}\index{R007@\texttt{R007} (rule)}
\rulefield{Message}\msgline{\meta{n} elements are both keep-in and keep-out}
\rulefield{Meaning}A keep-in and a keep-out region overlap. An element cannot be held solid and empty at once.
\rulefield{Fix}Remove the overlap: shrink one region, or split it.
\end{rulecard}

\begin{rulecard}{R008}{\sev{error}}{a support or load lies entirely inside a keep-out region}\label{rule:R008}\index{R008@\texttt{R008} (rule)}
\rulefield{Message}\msgline{\meta{support\textbar{}load} '\meta{id}' lies entirely inside a keep-out region}
\rulefield{Meaning}Every element around the item's nodes is held empty, so the load has no material path (or the support holds nothing). The solver can still converge on such a problem because empty elements keep a small stiffness; in the paper's gate ablation it converged and passed the numerical certificate.
\rulefield{Fix}Move the item out of the keep-out region, or move the keep-out region.
\end{rulecard}

\begin{rulecard}{R009}{\sev{warning}}{the minimum feature is not resolvable at this mesh}\label{rule:R009}\index{R009@\texttt{R009} (rule)}
\rulefield{Message}\msgline{min feature \meta{min feature} is \meta{n} elements; members below about \meta{2.4} elements are filtered out}
\rulefield{Meaning}\code{min\_feature} spans fewer than about 2.4 elements (twice the final filter radius of 1.2 elements); members that thin are filtered out.
\rulefield{Fix}Accept the one-click fix, which sets $h$ to \code{min\_feature}/2.4, or relax \code{min\_feature}. \emph{(The finding carries a fix: one click: set \code{resolution.h}.)}
\end{rulecard}

\begin{rulecard}{R010}{\sev{warning}}{a region was snapped to the grid}\label{rule:R010}\index{R010@\texttt{R010} (rule)}
\rulefield{Message}\msgline{\meta{n}}
\rulefield{Meaning}A patch had no node within half its size of its centre, or an edge fell between grid lines; it was moved to the nearest grid line. Nothing is snapped silently.
\rulefield{Fix}Usually harmless. If the change matters, align the coordinates with the grid or refine the mesh.
\end{rulecard}

\begin{rulecard}{R011}{\sev{warning}}{a passive region contains no element centre}\label{rule:R011}\index{R011@\texttt{R011} (rule)}
\rulefield{Message}\msgline{passive region '\meta{id}' contains no element centre at h=\meta{h}}
\rulefield{Meaning}At this element size no element centre falls inside the keep-in or keep-out region, so it has no effect.
\rulefield{Fix}Refine the mesh or enlarge the region. \emph{(The finding carries a fix: suggests: refine the mesh or enlarge the region.)}
\end{rulecard}

\begin{rulecard}{R012}{\sev{warning} / \sev{info}}{supports change on the coarsest multigrid level}\label{rule:R012}\index{R012@\texttt{R012} (rule)}
\rulefield{Message}\msgline{on the coarsest multigrid level (1/\meta{stride} resolution) the supports leave rigid modes free (\meta{modes}); the solver switches to the dilated coarse support rule (RobustGMG)} \newline \msgline{supports \meta{ids} are not represented on the coarsest multigrid level (1/\meta{stride}); the solve stays correct but may need more iterations}
\rulefield{Meaning}On the coarsest multigrid level the supports either leave rigid modes free (warning: the solver switches to the dilated coarse support rule) or lose some supports (info: the solve stays correct but may need more iterations). Neither blocks solving.
\rulefield{Fix}No action is needed. The finding explains a slower solve.
\end{rulecard}

\begin{rulecard}{R013}{\sev{error}}{the grid needs more GPU memory than is available}\label{rule:R013}\index{R013@\texttt{R013} (rule)}
\rulefield{Message}\msgline{\meta{g.n\_elem / 1000000.0}M elements need about \meta{need / 1000000000.0} GB with \meta{solver}; the GPU has \meta{vram\_bytes / 1000000000.0} GB}
\rulefield{Meaning}The memory estimate (bytes per element times elements) exceeds the GPU memory. The estimate uses about 366 bytes per element for \code{mg64} and 146 for \code{mg\_ir}.
\rulefield{Fix}Solve a coarser rung. A solve checks with \code{mg\_ir}; validation at the spec's own mesh checks with \code{mg64}. \emph{(The finding carries a fix: suggests: use mg\_ir or a coarser rung.)}
\end{rulecard}


\section{What the rules cannot see}\label{sec:rulelimits}
The rules decide whether a problem is admissible, not whether it is the problem the engineer meant. A support that
selects real nodes and removes every rigid-body mode passes, even if it restrains a direction the brief leaves free. A
keep-in pad that the brief never mentioned passes. A load with the right magnitude on the wrong face passes. These are
intent errors, and only the engineer, helped by clarification, can catch them (Chapter~\ref{ch:assistant}).

\begin{evidencebox}
With grounding tools but no questions, the primary agent produced 6 wrong specifications on the 50-brief suite. All 6
passed every rule \claim{C4.2}. Conversely, the rules stopped all 3 deliberately ill-posed briefs before any solve: two
rigid-body cases (\rref{R003}) and one load inside its own keep-out region (\rref{R008}) \claim{C4.1}.
\end{evidencebox}

The rules also do not check that the final design is connected, that the exported surface is manufacturable, or that
the structure is safe. The release certificate (Chapter~\ref{ch:ladder}) covers the numerical computation, and nothing
in AgenticTO covers structural safety.
